% Bagian: intuitionistic-logic
% Bab: tableaux
% Subbagian: introduction

\documentclass[../../../include/open-logic-section]{subfiles}

\begin{document}

\olfileid[id]{int}{tab}{int}

\olsection{Pendahuluan}

!!^{tableau}s adalah pohon tertentu (yang bercabang ke bawah) dari
!!{signed formula}, yakni pasangan yang terdiri atas tanda nilai kebenaran
($\True$ atau $\False$) dan !!a{sentence}
\[
\sFmla{\True}{!A} \text{ atau } \sFmla{\False}{!A}.
\]
!!^a{tableau} dimulai dengan sejumlah \emph{asumsi}. Setiap
!!{signed formula} berikutnya dihasilkan dengan menerapkan salah satu aturan
inferensi. Beberapa aturan inferensi menambahkan satu atau lebih
!!{signed formula} pada suatu ujung pohon; aturan lainnya menambahkan dua
ujung baru sehingga menghasilkan dua cabang. Aturan menghasilkan
!!{signed formula} yang !!{formula}-nya kurang kompleks daripada formula
dalam !!{signed formula} yang dikenai aturan tersebut. Jika suatu cabang
memuat sekaligus $\sFmla{\True}{!A}$ dan $\sFmla{\False}{!A}$, kita mengatakan
bahwa cabang itu \emph{tertutup}. Jika setiap cabang dalam !!a{tableau}
tertutup, seluruh !!{tableau} tersebut tertutup. !!^a{tableau} tertutup
merupakan !!a{derivation} yang menunjukkan bahwa himpunan !!{signed formula}
yang digunakan untuk memulai !!{tableau} tersebut tidak dapat dipenuhi. Hal
ini dapat digunakan untuk mendefinisikan relasi~$\Proves$:
$\Gamma \Proves !A$ jika dan hanya jika terdapat suatu himpunan berhingga
$\Gamma_0 = \{!B_1, \dots, !B_n\} \subseteq \Gamma$ sedemikian sehingga
terdapat !!{tableau} tertutup untuk asumsi-asumsi
\[
\{\sFmla{\False}{!A}, \sFmla{\True}{!B_1}, \dots, \sFmla{\True}{!B_n}\}.
\]

Untuk logika intuisionistik, kita harus memperluas pengertian
!!{signed formula} sekaligus menyesuaikan aturan-aturan bagi konektif. Selain
tanda ($\True$ atau $\False$), !!{formula} dalam !!{tableau} intuisionistik
juga mempunyai \emph{prefiks}~$\sigma$. Prefiks adalah barisan takkosong dari
bilangan bulat positif, yakni
$\sigma \in (\PosInt)^* \setminus \{\emptyseq\}$. Kita menuliskan prefiks
semacam itu tanpa tanda~$\tuple{\ }$ yang mengelilinginya, dan memisahkan
!!{element} individual dengan tanda~$.$ alih-alih tanda~$,$. Jika $\sigma$
adalah suatu prefiks, maka $\sigma.n$ adalah
$\sigma \concat \tuple{n}$; misalnya, jika $\sigma = 1.2.1$, maka
$\sigma.3$ adalah $1.2.1.3$. Jadi, sebagai contoh,
\[
\sFmla{\True}{!A \lif (!B \lif !C)}[1.2]
\]
adalah \emph{!!{signed formula} berprefiks} (atau, singkatnya,
\emph{!!{formula} berprefiks}).

Secara intuitif, prefiks menamai suatu dunia dalam model yang mungkin
memenuhi !!{formula} pada suatu cabang !!a{tableau}; dan jika $\sigma$
menamai suatu dunia, maka $\sigma.n$ menamai dunia yang terjangkau dari
(dunia yang dinamai oleh)~$\sigma$.

Dalam model intuisionistik, relasi keterjangkauan bersifat refleksif dan
transitif. Dalam istilah prefiks, hal ini berarti bahwa $\sigma$ terjangkau
dari $\sigma$ sendiri, begitu pula setiap prefiks yang memperpanjang
$\sigma$, yakni setiap prefiks berbentuk $\sigma.n_1.\cdots.n_k$. Mari kita
perkenalkan notasi $\sigma.*$ untuk menyatakan $\sigma$ sendiri beserta
setiap perluasannya. Dengan kata lain, prefiks-prefiks $\sigma.*$ adalah
semua dan hanya prefiks yang terjangkau dari~$\sigma$.

\end{document}
